%% kapa: exam=1 round=33 year=2022 period=1 form=A subject=real_estate subject_ko=부동산학원론 seq=5 number=85
%% source: Q-Net 감정평가사 기출문제 (cst00309) artlSeq=5209997 file=raw/1008/5209997/2022년도 제33회 감정평가사 1차시험 1교시(민법,경제학,부동산학).pdf
%% conversion: kordoc 4.24.51 -> tools/gen1.py (auto); review=proofread-sonnet; flags=none
\begin{question}{85}
디파스켈리-위튼(DiPasquale \& Wheaton)의 사분면 모형에 관한 설명으로 옳지 않은 것은? (단, 주어진 조건에 한함)
\begin{choices}
  \item 장기균형에서 4개의 내생변수, 즉 공간재고, 임대료, 자본환원율, 건물의 신규공급량이 결정된다.
  \item 신축을 통한 건물의 신규공급량은 부동산 자산가격, 생산요소가격 등에 의해 영향을 받는다.
  \item 자본환원율은 요구수익률을 의미하며 시장이자율 등에 의해 영향을 받는다.
  \item 최초 공간재고가 공간서비스에 대한 수요량과 일치할 때 균형임대료가 결정된다.
  \item 건물의 신규공급량과 기존 재고의 소멸에 의한 재고량 감소분이 일치할 때 장기균형에 도달한다.
\end{choices}
\end{question}
\answer{1}
